\documentclass[a4paper,12pt]{article}

\usepackage[ngerman]{babel}
\usepackage[T1]{fontenc}
\usepackage[utf8]{inputenc}
\usepackage{setspace}
\usepackage{parskip}

\title{Report}
\author{}
\date{}

\begin{document}

\maketitle

Ich würde die Streuung bei Namen anhand der Häufigkeit und Anzahl der unterschiedlichen Vor- und Nachnamen berechnen.
In der Realität kommen Namen unterschiedlich häufig vor, sie sind also nicht gleichmäßig verteilt.
Sollte die Faker-Bibliothek saubere zufällige Datensätze liefern (was unrealistisch wäre), müssten alle Vornamen bzw.
alle Nachnamen annähernd gleich oft vorkommen.

Die hohe Zufälligkeit der Streuung der erzeugten Namen würde sich anhand einer hohen Anzahl unterschiedlicher Namen zeigen,
die allerdings alle ähnlich häufig vorkommen. Das wäre realitätsfern.

Ich habe die Streuung der Bibliothek überprüft, indem ich folgende Daten ausgelesen habe:

\section*{Verteilung der Vornamen}
\begin{itemize}
    \item Anzahl unterschiedlicher Vornamen: 2002
    \item Häufigster Vorname: Jonas (492)
    \item Seltenster Vorname: Natalja (199)
    \item Durchschnittliche Häufigkeit: 249{,}75
\end{itemize}

\section*{Verteilung der Nachnamen}
\begin{itemize}
    \item Anzahl unterschiedlicher Nachnamen: 404
    \item Häufigster Nachname: Seifert (2533)
    \item Seltenster Nachname: Tintzmann (1123)
    \item Durchschnittliche Häufigkeit: 1237{,}62
\end{itemize}

Der häufigste Vorname kommt 2{,}47-mal so oft vor wie der seltenste.
Der häufigste Nachname kommt 2{,}26-mal so oft vor wie der seltenste.
Dies entspricht einer realistisch ungleichmäßigen Verteilung und ist somit realitätsnah.

\section*{Performance der Datenabfrage}

Zur Überprüfung der Performance wurde folgender SQL-Select verwendet:

\begin{quote}
\texttt{SELECT COUNT(DISTINCT vorname) FROM personen;}
\end{quote}

\subsection*{Ohne Index}
\begin{itemize}
    \item Run Time: real 0.305, user 0.281250, sys 0.031250
    \item Seitengröße: 4096B
    \item Seitenanzahl: 2870
    \item Speicherbedarf: 11{,}755520 MB
\end{itemize}

\subsection*{Mit Index}
\begin{itemize}
    \item Run Time: real 0.057, user 0.062500, sys 0.000000
    \item Seitengröße: 4096B
    \item Seitenanzahl: 6703
    \item Speicherbedarf: 27{,}455488 MB
\end{itemize}

Die Abfrage ist mit Index etwa fünfmal schneller, benötigt jedoch 2{,}34-mal mehr Speicher.

\section*{Vergleich mit biased Datenbank (50\% Vorname „Max“)}

\subsection*{Ohne Index}
\begin{itemize}
    \item Run Time: real 0.200, user 0.156250, sys 0.031250
    \item Seitengröße: 4096B
    \item Seitenanzahl: 2648
    \item Speicherbedarf: 10{,}846208 MB
\end{itemize}

\subsection*{Mit Index}
\begin{itemize}
    \item Run Time: real 0.064, user 0.046875, sys 0.000000
    \item Seitengröße: 4096B
    \item Seitenanzahl: 4336
    \item Speicherbedarf: 17{,}760256 MB
\end{itemize}

Die Performance ist etwa dreimal besser, der Speicherbedarf steigt um den Faktor 1{,}64.

Auffällig ist, dass der Select bei der indexierten biased Datenbank weniger Speicher benötigt und fast die gleiche Performance bietet wie die indexierte normale Datenbank. Dies liegt vermutlich daran, dass der Schlüsselwert „Vorname“ in 50\% der Fälle identisch ist und somit weniger Knoten im Index erzeugt werden müssen.

\end{document}
